\chapter{Manual de usuario}

\section{Interfaz principal: \textit{HUD}}
Al ejecutar la aplicación el usuario podrá observar dos elementos en la pantalla que serán una constante durante todas las operaciones.
En la parte superior se encuentran tanto el remarcador de marcos o \textit{frames} a los que van dirigidas las acciones 
realizadas en, por ejemplo, la planificación de rutas, como el botón de acceso al menú.

\section{Menú principal}
La aplicación cuenta con un menú simple a modo de selector.

Pulsando en el icono de tres puntos en la esquina superior derecha de la pantalla del dispositivo se desplegará o 
esconderá el menú. Desde este se podrán acceder pulsando 
las opciones pertinentes a las diferentes funcionalidades desarrolladas, además de activar o desactivar
características necesarias para su funcionamiento o emplear herramientas de depuración.

Las opciones disponibles son las siguientes:
\begin{itemize}
    \item Teleoperación: permite al usuario manejar robots móviles publicando las órdenes en un tema o \textit{topic} presente 
    en la interfaz (\textit{HUD}).
    \item Configuración de nube de puntos: permite al usuario configurar el tema (\textit{topic}) que publica la nuve de puntos
    del sensor en cada marco (\textit{frame}).
    \item Planificar ruta: permite al usuario planificar una ruta con origen en el marco seleccionado en el selector superior 
    del \textit{HUD}.
    \item Activar o desactivar \textit{frames}: permite al usuario visualizar u ocultar los marcos presentes en la sesión.
    \item Medición de frames: permite al usuario realizar mediciones entre los marcos presentes. Entre estas se incluye la distancia,
    la rotación, el \textit{pitch} y el \textit{yaw}.
    \item Ajuste de frames: permite a los usuarios realizar escalados de marcos (aumentar o disminuir su tamaño) y modificar su posición
    en el mundo virtual.
    \item \textit{Topic echo}: permite al usuario suscribirse a un tema y ver los contenidos publicados.
    \item \textit{Topic publisher}: permite al usuario publicar mensajes en un tema.
    \item Modificar parámetros: permite al usuario modificar los parámetros de los nodos presenetes.
    \item Configuración: permite al usuario modificar aspectos de la aplicación como el idioma, el sistema de puntos de referencia,
    la dirección IP del servidor ROS, etc.
    \item Visualización de superficies: permite al usuario visualizar u ocultar las superficies detectadas por el sistema de realidad
    aumentada.
    \item Menú de depuración de realidad aumentada: permite al usuario obtener información de depuración como los marcos/cuadros por 
    segundo a los que se está ejecutando la aplicación.
\end{itemize}

\section{Representación de \textit{frames} en la aplicación}
Los marcos ROS se representan en la aplicación como objetos flotantes esféricos de los que sobresalen sus ejes.

La esfera interna proporciona información de estado de la actualización de las transformaciones, de manera que 
el color verde representará que no existe ningún problema, el color azul indicará que se está en modo pausa y el color rojo
indicará que el marco no ha sido refrescado en un periodo de tiempo elevado.

Las relaciones entre marcos (hijos y padres) se visualizarán como una línea roja con origen y destino de marcos diferentes.

En los marcos móviles es posible también visualizar el vector de velocidad como una línea azul que sale 
del centro de la esfera en la que la longitud indica su magnitud.

Es posible, además, visualizar el \textit{trial} de las últimas 10 posiciones de la tranformación como un vector opuesto al de
la velocidad que tiene origen, de igual manera, en el centro de la esfera. De esta forma
se indica si hay problemas de conexión o retardos en la red mediante un código de colores formado por rojo, verde y amarillo.

\section{Remarcador de marcos}
El remarcador o \textit{highlighter} de marcos se encuentra en la parte superior de la pantalla durante toda su ejecución.
Este sirve para diferentes propósitos.

En primer lugar, permite reconfigurar las etiquetas de detección (\textit{AprilTags}) de forma que si el usuario comete un 
error al asignar un marco o \textit{frame} a una etiqueta, puede pulsar en el icono de la papelera y esta asociación quedará 
disuelta.
Una vez que la cámara detecte de nuevo la etiqueta liberada, volverá a solicitar al usuario la configuración.

Por otra parte, permite que las acciones realizadas en ciertas funcionalidades tengan como objetivo el marco mostrado en el 
seleccionados. Por ejemplo, en la planificación de ruta, el objetivo del movimiento será el \textit{frame} seleccionado en este 
componente. De hecho, esto se indicará de forma visual al usuario como una linea azul que saldrá del marco hasta la primera 
posición de la ruta.

\section{Teleoperación}
Para operar telemáticamente a un robot, se selecciona la opción "Teleop" en el menú de la aplicación.
Esta interfaz añade diferentes elementos al \textit{HUD}.

En primer lugar, en el lateral izquierdo de la pantalla se encuentra un botón que al ser pulsado despliega una 
ventana. En esta se pueden modificar los parámetros del movimiento como son la velocidad angular y lineal.

En la parte inferior izquierda se encuentran dos circulos concénctricos de diferente color. Estos se corresponden al 
\textit{joystick} que el usuario emplea para enviar las órdenes de movimiento en el plano horizontal (moverse hacia adelante,
atrás, izquierda o derecha).

En la pate inferior central, se encuentra un selector en el que se muestra el tema (\textit{topic}) objetivo del movimiento,
siendo necesario que el usuario seleccione una entre las opciones disponibles.

En la parte inferior derecha se encuentran dos botones con flechas hacia arriba y abajo. Al pulsarlos, el usuario podrá 
enviar la órden de movimiento en el plano vertical.

\section{Configuración de nube de puntos}
Para realizar la configuración de la nube de puntos se selecciona "Configuración de Nube de Puntos" en el menú.
La aplicación muestra la lista de los \textit{frames} en el lado izquierdo y un selector en el lado derecho.

Una vez realizada la selección, la aplicación almacena automáticamente la información y procede a mostrar la nube en el marco
correspondiente.

El usuario puede cerrar la ventana pulsando en el botón de cierre en la esquina superior derecha sin necesidad de guardar los cambios.

\section{Planificación de rutas}
Para realizar la planificación de rutas se selecciona "Path planning" en el menú de la aplicación.
Esta interfaz añade diferentes elementos al \textit{HUD}.
En la parte inferior izquierda se encuentra el botón de deshacer. Este elimina el último punto de ruta introducido.

En la parte inferior central se encuentra el selector del tema (\textit{topic}) objetivo del movimiento (dónde se publicarán los mensajes).

En la parte inferior derecha se encuentra el botón de ejecución de ruta. Este envía punto por punto la información 
al tema seleccionado, lo que tiene como resultado que el robot se mueva siguiendo la ruta introducida.

Para esta funcionalidad es necesario tener la visualización de superficies capturadas por la realidad aumentada activada. 
La aplicación lo realizará de manera automática y almacenará el estado anterior para retomarlo al abandonar la característica.
Sin embargo, el usuario puede en todo momento desactivarlas dentro de esta funcionalidad, pero no podrá situar puntos de ruta nuevos.

\section{Medición de \textit{frames}}
Para realizar mediciones entre marcos se selecciona en el menú la opción "Medición de frames".
Esta interfaz introduce una nueva ventana al \textit{HUD} de la aplicación.
En esta se indica la distancia y rotación de los marcos seleccionados.

Para seleccionar un marco el usuario pulsará durante unos instantes a su representación en realidad aumentada. Tras esto, 
la representación cambiará de color y quedará seleccionada. 

Seleccionado dos marcos, la información se actualiza en la ventana.

\section{Ajuste de marcos}
Para realizar ajustes en la posición de las representación de los marcos o modificar su tamaño se selecciona "Ajuste de frames" en 
el menú.

Esta funcionalidad añade diferentes elementos en el \textit{HUD} de la aplicación.

El botón de pausa se puede emplear para detener la actualización de las transformaciones, de manera que se genera una 
instantánea de la situación.

El resto de elementos permiten ajustar la posición del marco en el mundo. 

Pulsando en los botones de color rojo la representación del marco se moverá en el eje X. 
Pulsando en los botones de color azul la representación del marco se moverá en el eje Z.
Pulsando en los botones de color verde la representación del marco se moverá en el eje Y.

\section{\textit{Topic echo}}
Para visualizar los mensajes publicados en un tema se selecciona en el menú la opción "Topic echo".

Esta interfaz cuenta con dos selectores y una ventana de visualización.

El selector de tema permite escoger entre las diferentes opciones el \textit{topic} deseado.
El segundo selector permite seleccionar el tipo de mensaje deseado del tema (típicamente solo se dispone de una opción).

Al pulsar en iniciar, el contenido publicado en el tema se muestra en la ventana de visualización y los selectores 
se bloquean hasta que se pulsa en parar.

\section{\textit{Topic Publisher}}
Para publicar un mensaje en un tema, se selecciona en el menú la opción "Topic Publisher".

El selector de temas muestra las opciones disponibles (los \textit{topics} presentes en la red ROS).
Al seleccionar una opción se modificará también el tipo de mensaje mostrado en el selector siguiente y la plantilla del mensaje.
Generalmente solo existe un tipo de mensaje por tema, pero se ha habilitado la selección pora posibles
modificaciones futuras.

Se introduce el mensaje deseado y se modifica la frecuencia de publicación.

Al pulsar en aceptar se publica en el tema el mansaje correspondiente siguiendo la frecuencia.

\section{Modificar parámetros}
Para modificar un paŕametro en la ventana de modificación de parámetros, se selecciona un nodo entre los presentes,
el parámetro deseado, se introduce el nuevo valor en la caja de texto inferior y se pulsa el botón de guardado.

Al guardar el parámetro tomará el nuevo valor y se mostrará refrescado.

\section{Configuración}
Para acceder a la configuración, se selecciona en el menú la opción "Configuración".

La configuración permite modificar el idioma, seleccionar el sistema de puntos de referencia empleados en la aplicación y 
variar parámetros de conexión con el servidor.

Al itroducir cualquier modificación, se debe pulsar en el botón de guardado, de lo contrario los cambios no se verán reflejados.

\section{Visualización de superficies}
Para visualizar u ocultar las superficies capturadas por el núcleo de realidad aumentada se selecciona "Visualizar superficies"
en el menú.

Esta opción es un botón con estado, es decir, que se indicará si se encuentra activo o inactivo en todo momento, de manera que 
al estar el deslizador en la parte derecha de color verde, las superficies se visualizarán.

De lo contrario, el botón estará de color gris y las superficies serán desactivadas.

\section{Menú de depuración}
Para visualizar u ocultar el menú de depuración de realidad aumentada, se seleccionma la opción "AR Debug Menu".